\chapter{Делимость, простые числа, НОД и НОК}
\section*{Делимость и кратность}

\begin{definition}
\term{Делимость нацело} (\emph{divisibility}): целое число \(b\) делится на целое число \(a\ne0\),
если существует целое число \(k\), для которого
\[
b=ak.
\]
При этом \(a\) называют \term{делителем} (\emph{divisor}) числа \(b\).
\end{definition}

Запись с тремя вертикальными точками \(b\mathrel{\vdots}a\) читается:
«\(b\) делится на \(a\) нацело».
В университетских курсах также используют запись \(a\mid b\): «\(a\) делит \(b\)».
Порядок чисел в этих двух обозначениях противоположен:
\[
12\mathrel{\vdots}4\quad\Longleftrightarrow\quad4\mid12,
\qquad 12=4\cdot3.
\]
Здесь знак \(\Longleftrightarrow\) означает, что записи выражают одно и то же утверждение.

\begin{definition}
\term{Кратное} (\emph{multiple}) числа \(a\ne0\) --- целое число \(b\), которое можно представить как \(b=ak\)
с целым \(k\). Иначе говоря, \(b\) кратно \(a\), если \(b\mathrel{\vdots}a\).
\end{definition}

Например, положительные кратные числа \(4\): \(4,8,12,16,\ldots\).
Ноль также кратен любому ненулевому целому числу: \(0=a\cdot0\).
В задачах на НОК рассматриваются только \emph{положительные} общие кратные.

\section*{Простые и составные числа}

\begin{definition}
\term{Простое число} (\emph{prime number}) --- натуральное число, имеющее ровно два положительных делителя:
\(1\) и само это число.
\end{definition}

\begin{definition}
\term{Составное число} (\emph{composite number}) --- натуральное число больше \(1\), имеющее положительный делитель,
отличный от \(1\) и от самого числа.
\end{definition}

Число \(7\) простое: его положительные делители --- \(1\) и \(7\).
Число \(9\) составное: кроме \(1\) и \(9\), оно делится на \(3\).
Число \(1\) не является ни простым, ни составным.
Число \(2\) --- единственное чётное простое число.

\section*{Наибольший общий делитель и наименьшее общее кратное}

\begin{definition}
\term{Наибольший общий делитель} (\emph{greatest common divisor, gcd}) положительных целых чисел \(a\) и \(b\) ---
наибольшее положительное целое число, на которое оба числа делятся нацело.
Обозначения: \(\text{НОД}(a,b)\) и \(\gcd(a,b)\).
Английское название --- \emph{greatest common divisor}, сокращённо \emph{gcd}.
\end{definition}

Общие положительные делители чисел \(12\) и \(18\) --- \(1,2,3,6\), поэтому
\[
\text{НОД}(12,18)=\gcd(12,18)=6.
\]

\begin{definition}
\term{Наименьшее общее кратное} (\emph{least common multiple, lcm}) положительных целых чисел \(a\) и \(b\) ---
наименьшее положительное целое число, которое делится нацело и на \(a\), и на \(b\).
Обозначения: \(\text{НОК}(a,b)\) и \(\operatorname{lcm}(a,b)\).
Английское название --- \emph{least common multiple}, сокращённо \emph{lcm}.
\end{definition}

Положительные общие кратные чисел \(12\) и \(18\) --- \(36,72,108,\ldots\), поэтому
\[
\text{НОК}(12,18)=\operatorname{lcm}(12,18)=36.
\]

\begin{definition}
\term{Взаимно простые числа} (\emph{coprime integers}) --- положительные целые числа, НОД которых равен \(1\).
\end{definition}
Например, \(8\) и \(15\) взаимно просты, хотя оба числа составные.

\section*{Основная теорема арифметики (ОТА)}

\begin{theorem}
Каждое натуральное число больше \(1\) можно представить произведением простых чисел.
Такое разложение единственно с точностью до перестановки множителей.
\end{theorem}
\begin{proof}Существование докажем индукцией по числу: простое число уже разложено; составное имеет вид \(ab\), где \(1<a,b<n\), и к обоим множителям применяется предположение индукции. Для единственности используем следующий факт: если простое \(p\) делит \(ab\), то оно делит \(a\) или \(b\). Если \(p\nmid a\), то \(\gcd(p,a)=1\); обратная подстановка в алгоритме Евклида даёт целые \(u,v\) с \(up+va=1\). Умножение на \(b\) показывает, что \(p\mid b\). Теперь первый простой множитель одного разложения обязан совпасть с одним из множителей другого; сократим его и повторим. Все множители совпадут с точностью до порядка.\end{proof}

Например,
\[
60=2\cdot2\cdot3\cdot5=2^2\cdot3\cdot5.
\]
Множители \(2,3,5\) простые; запись \(60=4\cdot15\) ещё не является разложением на простые множители.

\section*{Вычисление НОД и НОК по разложению}

Разложите каждое число на простые множители, объединив повторения в степени.
Для НОД возьмите только общие простые множители в \emph{меньших} степенях.
Для НОК возьмите все встречающиеся простые множители в \emph{больших} степенях.
Если общих простых множителей нет, НОД равен \(1\).

Например,
\[
12=2^2\cdot3,\qquad18=2\cdot3^2,
\]
\[
\gcd(12,18)=2\cdot3=6,\qquad
\operatorname{lcm}(12,18)=2^2\cdot3^2=36.
\]
Для нескольких положительных чисел правило такое же:
в НОД входят множители, общие \emph{всем} числам, с наименьшими показателями;
в НОК --- все встретившиеся множители с наибольшими показателями.

\begin{theorem}
Для любых положительных целых чисел \(a,b\)
\[
\gcd(a,b)\operatorname{lcm}(a,b)=ab.
\]
\end{theorem}
\begin{proof}Для каждого простого множителя обозначим его показатели в \(a,b\) через \(\alpha,\beta\), допуская показатель \(0\). В НОД его показатель \(\min(\alpha,\beta)\), в НОК — \(\max(\alpha,\beta)\). Их сумма равна \(\alpha+\beta\), то есть показателю того же простого в \(ab\). По единственности разложения произведения равны.\end{proof}
Это равенство относится к двум числам; напрямую переносить его на три числа нельзя.

\section*{Алгоритм Евклида}

Для нахождения НОД можно последовательно делить с остатком,
каждый раз заменяя пару «делимое, делитель» парой «делитель, остаток».
Когда остаток стал нулём, НОД равен последнему ненулевому делителю.
Например,
\[
48=2\cdot18+12,\qquad
18=1\cdot12+6,\qquad
12=2\cdot6+0,
\]
поэтому \(\gcd(48,18)=6\).

\section*{Признаки делимости}

В десятичной записи натурального числа действуют следующие признаки.
Каждый из них даёт необходимое и достаточное условие: число делится на указанный делитель
\emph{тогда и только тогда}, когда выполнена соответствующая проверка.

\begin{itemize}[leftmargin=1.5em,itemsep=4pt]
\item \textbf{На \(2\):} последняя цифра --- \(0,2,4,6\) или \(8\).
\item \textbf{На \(3\):} сумма цифр делится на \(3\).
\item \textbf{На \(4\):} число, образованное двумя последними цифрами, делится на \(4\).
\item \textbf{На \(5\):} последняя цифра --- \(0\) или \(5\).
\item \textbf{На \(6\):} число делится одновременно на \(2\) и на \(3\).
\item \textbf{На \(8\):} число, образованное тремя последними цифрами, делится на \(8\).
\item \textbf{На \(9\):} сумма цифр делится на \(9\).
\item \textbf{На \(10\):} последняя цифра --- \(0\).
\end{itemize}

Если число содержит меньше двух или трёх цифр, в признаках на \(4\) и \(8\) проверяют само число.
Например, \(123\mathrel{\vdots}3\), поскольку \(1+2+3=6\),
но \(123\) не делится на \(9\), поскольку \(6\) не делится на \(9\).
\Needspace{7\baselineskip}\section*{Задачи}
\begin{enumerate}[label={},leftmargin=3.2em,itemsep=10pt]
\Needspace{3\baselineskip}\item[\textbf{2.1.}] Выпиши все положительные делители \(36\).
\Needspace{3\baselineskip}\item[\textbf{2.2.}] Выпиши все положительные делители \(48\).
\Needspace{3\baselineskip}\item[\textbf{2.3.}] Найди общие положительные делители \(20\) и \(30\).
\Needspace{3\baselineskip}\item[\textbf{2.4.}] Выпиши кратные \(7\) между \(20\) и \(60\).
\Needspace{3\baselineskip}\item[\textbf{2.5.}] Найди наименьшее трёхзначное число, кратное \(8\).
\Needspace{3\baselineskip}\item[\textbf{2.6.}] Раздели на простые и составные: \(17,\allowbreak 25,\allowbreak 31,\allowbreak 39,\allowbreak 49,\allowbreak 51,\allowbreak 53\).
\Needspace{3\baselineskip}\item[\textbf{2.7.}] Выпиши простые числа от \(20\) до \(40\).
\Needspace{3\baselineskip}\item[\textbf{2.8.}] Может ли простое число оканчиваться на \(0\)? Объясни.
\Needspace{3\baselineskip}\item[\textbf{2.9.}] Любое ли нечётное число простое? Приведи контрпример.
\Needspace{3\baselineskip}\item[\textbf{2.10.}] Какие числа делятся на \(3\), а какие на \(9\): \(126,\allowbreak 234,\allowbreak 405,\allowbreak 517,\allowbreak 702\)?
\Needspace{3\baselineskip}\item[\textbf{2.11.}] Какие числа делятся на \(2\) и на \(5\) одновременно: \(125,\allowbreak 230,\allowbreak 342,\allowbreak 480,\allowbreak 705\)?
\Needspace{3\baselineskip}\item[\textbf{2.12.}] Какие цифры можно поставить вместо \(\square \) в \(43 \square \), чтобы число делилось на \(3\)?
\Needspace{3\baselineskip}\item[\textbf{2.13.}] Какая цифра вместо \(\square \) в \(7 \square  2\) даёт делимость на \(9\)?
\Needspace{3\baselineskip}\item[\textbf{2.14.}] Какие цифры вместо \(\square \) в \(12 \square \) дают делимость на \(2\) и на \(3\)?
\Needspace{3\baselineskip}\item[\textbf{2.15.}] Разложи \(72\) на простые множители.
\Needspace{3\baselineskip}\item[\textbf{2.16.}] Разложи \(150\) на простые множители.
\Needspace{3\baselineskip}\item[\textbf{2.17.}] Разложи \(252\) на простые множители.
\Needspace{3\baselineskip}\item[\textbf{2.18.}] Заверши разложение \(90 = 2 \cdot  5 \cdot  9\) на простые множители.
\Needspace{3\baselineskip}\item[\textbf{2.19.}] Найди НОД и НОК чисел \(24\) и \(36\).
\Needspace{3\baselineskip}\item[\textbf{2.20.}] Найди НОД и НОК чисел \(18\) и \(45\).
\Needspace{3\baselineskip}\item[\textbf{2.21.}] Найди НОД и НОК чисел \(8\) и \(15\).
\Needspace{3\baselineskip}\item[\textbf{2.22.}] Найди НОД и НОК чисел \(12,\allowbreak 18\) и \(30\).
\Needspace{3\baselineskip}\item[\textbf{2.23.}] Из \(24\) карандашей и \(36\) ручек собери максимум одинаковых наборов, использовав всё. Сколько наборов и что в каждом?
\Needspace{3\baselineskip}\item[\textbf{2.24.}] Автобусы отправились вместе в 9:00. Интервалы отправления: \(12\) и \(18\) минут. Когда следующее совместное отправление?
\Needspace{3\baselineskip}\item[\textbf{2.25.}] Прямоугольник \(60 \times  84\) см разрежь без остатка на одинаковые квадраты наибольшей стороны. Найди сторону и число квадратов.
\Needspace{3\baselineskip}\item[\textbf{2.26.}] Выпиши все положительные делители \(24\).
\Needspace{3\baselineskip}\item[\textbf{2.27.}] Выпиши первые пять положительных кратных \(6\).
\Needspace{3\baselineskip}\item[\textbf{2.28.}] Какие из чисел \(1,\allowbreak 2,\allowbreak 9,\allowbreak 13,\allowbreak 21,\allowbreak 29\) простые?
\Needspace{3\baselineskip}\item[\textbf{2.29.}] На какие числа из списка \(2,\allowbreak 3,\allowbreak 5,\allowbreak 9,\allowbreak 10\) делится \(450\)?
\Needspace{3\baselineskip}\item[\textbf{2.30.}] Разложи \(84\) на простые множители.
\Needspace{3\baselineskip}\item[\textbf{2.31.}] Найди НОД и НОК чисел \(12\) и \(18\).
\Needspace{3\baselineskip}\item[\textbf{2.32.}] Из жизни. Для праздника есть \(42\) яблока и \(56\) мандаринов. Их нужно разложить без остатка в максимальное число одинаковых подарков. Сколько подарков получится и что будет в каждом?
\Needspace{3\baselineskip}\item[\textbf{2.33.}] Из жизни. Два индикатора загорелись одновременно. Один загорается каждые \(8\) секунд, другой каждые \(12\) секунд. Через сколько секунд они впервые снова загорятся вместе?
\Needspace{3\baselineskip}\item[\textbf{2.34.}] Из жизни. В коробке \(97\) деталей. Можно ли разложить их в одинаковые группы больше одной детали, получив больше одной группы, без остатка?
\Needspace{3\baselineskip}\item[\textbf{2.35.}\starmark] Найди наименьшее натуральное число, которое при делении на \(6,\allowbreak 8\) и \(9\) даёт остаток \(1\) и при этом больше \(1\).
\Needspace{3\baselineskip}\item[\textbf{2.36.}\starmark] НОД двух положительных целых чисел равен \(6\), НОК равен \(180\). Одно число \(30\). Найди второе.
\Needspace{3\baselineskip}\item[\textbf{2.37.}\starmark] Найди все цифры \(a\) и \(b\), чтобы четырёхзначное число \(\overline{3a5b}\) делилось на \(5\) и на \(9\).
\Needspace{3\baselineskip}\item[\textbf{2.38.}] Пусть \(a>b>0\) --- целые числа. Докажи равенство \(\gcd(a,b)=\gcd(b,a-b)\), показав, что пары имеют одни и те же положительные общие делители.
\Needspace{3\baselineskip}\item[\textbf{2.39.}] Пусть \(p\) --- простое число, \(a,b\) --- положительные целые числа. Используя основную теорему арифметики, докажи: если \(p\mid ab\), то \(p\mid a\) или \(p\mid b\).
\end{enumerate}
